\documentclass{article}

\usepackage{bm}
\usepackage{ctex}
\usepackage{tikz}
\usepackage{float}
\usepackage{xcolor}
\usepackage{setspace}
\usepackage{datetime}
\usepackage{geometry}
\usepackage{graphicx}
\usepackage{enumitem}
\usepackage{tcolorbox}
\usepackage{nicematrix}
\usepackage{amsmath, amssymb, amsfonts, mathrsfs}
\usepackage[fontsize = 12pt]{fontsize}

\renewcommand{\thesubsection}{(\alph{subsection})}
\newcommand{\diff}[2]{\dfrac{\mathrm{d}#1}{\mathrm{d}#2}}
\newcommand{\diffn}[3]{\dfrac{\mathrm{d}^{#3}#1}{\mathrm{d}#2^{#3}}}
\newcommand{\piff}[2]{\dfrac{\partial{#1}}{\partial{#2}}}
\newcommand{\eval}[2]{\left.#1\right|_{#2}}
\newcommand{\e}{\mathrm{e}}
\renewcommand{\d}{\mathrm{d}}
\newcommand{\barline}[1]{\overline{#1\rule{0pt}{1.9ex}}}

\geometry{
    a4paper,
    left = 3cm,
    right = 3cm,
    top = 2.4cm,
    bottom = 2.4cm
}

\setitemize[1]{
    itemsep = 0pt,
    partopsep = 0pt,
    parsep = \parskip,
    topsep = 5pt
}

\title{应用统计：第二次作业}
\author{Illusionna \quad 2025...............004 \quad 人工智能学院}
\date{\currenttime,\ \today}



\begin{document}

\maketitle



\section{Q-1.6}
\textit{\textcolor{gray}{两台车床加工同一种零件，第一台车床出现废品的概率为 0.03，第二台出现废品的概率为 0.02。若两台车床加工的零件放在一起，并且已知第一台车床加工的零件比第二台车床加工的零件多一倍。}}

设 $A$ 表示零件是第一台车床加工，$B$ 表示零件是第二台车床加工，$C$ 表示零件是废品，$\barline{C}$ 表示零件是合格品。
\[ \Pr(A)=\dfrac{2}{3}\qquad\Pr(B)=\dfrac{1}{3} \]

两台车床出现废品的概率：
\[ \Pr(C\mid A)=0.03\qquad\Pr(C\mid B)=0.02 \]



\subsection{\textit{\textcolor{gray}{求从整批零件中任取 1 件，恰好是合格的概率}}}
由全概率公式，任取 1 件是废品的概率为：
\[ \Pr(C)=\Pr(A)\Pr(C\mid A)+\Pr(B)\Pr(C\mid B)=\dfrac{2}{3}\times0.03+\dfrac{1}{3}\times0.02=\dfrac{2}{75} \]

因此，任取 1 件恰好是合格品的概率为：
\[ \Pr(\barline{C})=1-\Pr(C)=\dfrac{73}{75}\approx97.33\% \]



\subsection{\textit{\textcolor{gray}{如果任取 1 件后，发现是废品，求它是第二台车床加工的概率}}}
\[ \Pr(B\mid C)=\dfrac{\Pr(B)\Pr(C\mid B)}{\Pr(C)}=\dfrac{\dfrac{1}{3}\times0.02}{\displaystyle\dfrac{2}{75}}=25\% \]



\section{Q-1.8}
\textit{\textcolor{gray}{设离散型随机变量 $X$ 的概率分布为}}

\begin{table}[H]
    \centering
    \renewcommand{\arraystretch}{1.2}
    \setlength{\tabcolsep}{16pt}
    \color{gray}{
        \caption{随机变量 $X$ 的分布列}
        \begin{tabular}{|c|c|c|c|c|c|c|}
            \hline
            $X$ & $-2$ & $-1$ & $0$ & $1$ & $2$ & $3$ \\
            \hline
            $\Pr$ & 0.05 & 0.15 & 0.20 & 0.25 & 0.2 & 0.15 \\
            \hline
        \end{tabular}
    }
\end{table}

\textit{\textcolor{gray}{求 $Y=2X+1$ 和 $Z=X^2$ 的概率分布}}

由于所有的 $Y$ 值互不相同，对应的概率不变，故 $Y=2X+1$ 的概率分布为：

\begin{table}[H]
    \centering
    \renewcommand{\arraystretch}{1.2}
    \setlength{\tabcolsep}{16pt}
    \caption{随机变量 $Y$ 的分布列}
    \begin{tabular}{|c|c|c|c|c|c|c|}
        \hline
        $Y$ & $-3$ & $-1$ & $1$ & $3$ & $5$ & $7$ \\
        \hline
        $\Pr$ & 0.05 & 0.15 & 0.20 & 0.25 & 0.2 & 0.15 \\
        \hline
    \end{tabular}
\end{table}

我们计算 $X$ 取不同值时对应的 $Z$ 值，并将产生相同 $Z$ 值的概率相加：

\begin{table}[H]
    \centering
    \renewcommand{\arraystretch}{1.2}
    \setlength{\tabcolsep}{16pt}
    \caption{随机变量 $Z$ 的分布列}
    \begin{tabular}{|c|c|c|c|c|}
        \hline
        $Z$ & $0$ & $1$ & $4$ & $9$ \\
        \hline
        $\Pr$ & 0.20 & 0.40 & 0.25 & 0.15 \\
        \hline
    \end{tabular}
\end{table}



\section{Q-1.12}
\textit{\textcolor{gray}{设随机变量 $X\sim N(3,4)$，求概率 $\Pr(X\geqslant4)$ 及 $\Pr(-2\leqslant X\leqslant7)$。（$\Phi(2)=0.9772,\ \Phi(2.5)=0.9938,\ \Phi(0.5)=0.6915$）}}

随机变量 $X\sim N(3,4)$，标准化后 $Z=\dfrac{X-\mu}{\sigma}=\dfrac{X-3}{2}\sim N(0,1)$，则：
\begin{align*}
    \Pr(X\geqslant4)&=\Pr\left(\dfrac{X-3}{2}\geqslant\dfrac{4-3}{2}\right)
    \\
    &=\Pr(Z\geqslant0.5)
    \\
    &=1-\Phi(0.5)
    \\
    &=0.3085
\end{align*}
\begin{align*}
    \Pr(-2\leqslant X\leqslant7)&=\Pr\left(\dfrac{-2-3}{2}\leqslant\dfrac{X-3}{2}\leqslant\dfrac{7-3}{2}\right)
    \\
    &=\Pr(-2.5\leqslant Z\leqslant2)
    \\
    &=\Phi(2)-\Phi(-2.5)=\Phi(2)-\left[1-\Phi(2.5)\right]
    \\
    &=0.9772-(1-0.9938)=0.9710
\end{align*}



\end{document}